\documentclass[11pt,a4paper]{article}
\usepackage[utf8]{inputenc}
\usepackage[margin=1in]{geometry}
\usepackage{amsmath}
\usepackage{amssymb}
\usepackage{enumitem}
\usepackage{booktabs}
\usepackage{titlesec}
\usepackage{microtype}

\titleformat{\section}{\large\bfseries\sffamily}{}{0em}{}[\titlerule]
\titleformat{\subsection}{\normalsize\bfseries\sffamily}{}{0em}{}

\setlength{\parindent}{0pt}
\setlength{\parskip}{6pt plus 2pt minus 1pt}

\title{\textbf{Hard CAMS Exam Questions: Missing Topics Edition II}}
\author{\textbf{Advanced AML Training Framework}}
\date{\today}

\begin{document}
	
	\maketitle
	
	% -------------------------------------------------------
	\section{Customer Due Diligence (CDD) / KYC Program Pillars}
	% -------------------------------------------------------
	
	\subsection{1. The Four Pillars of a BSA/AML Compliance Program}
	Under U.S. Bank Secrecy Act (BSA) requirements (31 CFR \S 1020.210) as amended by FinCEN's 2016 CDD Final Rule, which of the following correctly identifies the \textit{five} minimum required pillars of a financial institution's AML compliance program?
	\begin{enumerate}[label=\Alph*.]
		\item Internal controls; a designated compliance officer; ongoing employee training; independent testing; customer due diligence (including beneficial ownership identification).
		\item Annual board resolutions; a quarterly sanctions report; a law enforcement liaison desk; a digital asset custodian; automated wire screening.
		\item Simplified due diligence for all customers; waiver of PEP screening for domestic clients; outsourced monitoring to a third-party vendor; a customer ombudsman; annual profit-sharing with regulators.
		\item Internal controls; a designated compliance officer; outsourced customer onboarding; independent testing; monthly currency transaction reports.
	\end{enumerate}
	\textbf{Correct Answer: A}
	
	\subsection{2. FinCEN's CDD Final Rule: Legal Entity Customers}
	Under FinCEN's 2016 CDD Final Rule, covered financial institutions must identify and verify the identity of any natural person who directly or indirectly owns \textbf{25\% or more} of a legal entity customer, \textit{and} at least one individual who exercises significant managerial control. What is the term for this second category of individual?
	\begin{enumerate}[label=\Alph*.]
		\item Ultimate Beneficial Owner
		\item Control Prong / Controlling Person
		\item Nominee Director
		\item Settlor of Record
	\end{enumerate}
	\textbf{Correct Answer: B}
	
	\subsection{3. Simplified Due Diligence (SDD): Eligibility}
	A compliance officer is reviewing a new account application from a publicly listed company whose shares are traded on the New York Stock Exchange. The company is subject to full SEC disclosure and regulatory reporting obligations. Under a risk-based approach, which CDD treatment is most appropriate?
	\begin{enumerate}[label=\Alph*.]
		\item Enhanced Due Diligence because all corporate entities are high risk.
		\item Simplified Due Diligence, as the entity is subject to robust regulatory transparency and disclosure requirements that mitigate money laundering risk.
		\item No due diligence because the NYSE listing acts as an absolute exemption.
		\item Standard CDD with mandatory source-of-funds verification for every transaction.
	\end{enumerate}
	\textbf{Correct Answer: B}
	
	\subsection{4. Ongoing Monitoring Obligations}
	After initial onboarding, a bank's CDD obligations require which of the following on a continuing basis?
	\begin{enumerate}[label=\Alph*.]
		\item A one-time verification of identity at account opening is sufficient; no further monitoring is needed unless the customer complains.
		\item Ongoing monitoring of account activity to identify transactions inconsistent with the customer's risk profile, and periodic refresh of customer information to keep records current and accurate.
		\item Annual notarized resubmission of all original identification documents regardless of risk level.
		\item Continuous real-time video surveillance of every customer's physical premises.
	\end{enumerate}
	\textbf{Correct Answer: B}
	
	% -------------------------------------------------------
	\section{SAR Tipping-Off, Confidentiality, and Safe Harbor}
	% -------------------------------------------------------
	
	\subsection{5. SAR Safe Harbor Protections}
	A bank files a Suspicious Activity Report on a long-standing corporate client. Three months later, the client sues the bank for defamation, claiming the SAR filing damaged their reputation. Under U.S. law (31 U.S.C. \S 5318(g)(3)), what protects the bank from civil liability?
	\begin{enumerate}[label=\Alph*.]
		\item The bank is not protected; it must prove the SAR was accurate in court.
		\item The statutory SAR Safe Harbor provision, which grants immunity from civil liability to any institution that files an SAR in good faith.
		\item The bank is only protected if the SAR resulted in a criminal conviction.
		\item The bank must obtain prior judicial approval before the Safe Harbor applies.
	\end{enumerate}
	\textbf{Correct Answer: B}
	
	\subsection{6. Mandatory vs. Voluntary SAR Filing}
	An AML analyst identifies a series of transactions that are suspicious but the total aggregate value is only \$4,000 -- below the standard \$5,000 mandatory filing threshold for U.S. bank SARs. Which statement is most accurate regarding the institution's obligations?
	\begin{enumerate}[label=\Alph*.]
		\item No obligation exists whatsoever; threshold-based rules are absolute floors with no discretion.
		\item The institution may still file a voluntary SAR if it believes the activity is suspicious, and in some cases (e.g., insider abuse) lower thresholds apply.
		\item The institution must wait until the aggregate reaches exactly \$5,000 on the same calendar day before any filing is permissible.
		\item Filing below the threshold automatically triggers a regulatory penalty.
	\end{enumerate}
	\textbf{Correct Answer: B}
	
	\subsection{7. Tipping-Off in a Multi-Jurisdictional Investigation}
	A compliance officer at a European bank shares the existence of a pending STR with a colleague at a subsidiary bank in another country over an unencrypted email, because both institutions serve the same suspect client. Why is this action problematic under AMLD and most national AML frameworks?
	\begin{enumerate}[label=\Alph*.]
		\item Inter-group sharing of STR information is always fully prohibited regardless of circumstances.
		\item While limited group-level sharing may be permissible under strict conditions, uncontrolled disclosure of an STR to parties outside the tipping-off exemption risks alerting the subject and violates the tipping-off prohibition.
		\item The action is permissible because both entities share the same client.
		\item The only prohibition is against sharing with law enforcement.
	\end{enumerate}
	\textbf{Correct Answer: B}
	
	% -------------------------------------------------------
	\section{Real Estate Money Laundering}
	% -------------------------------------------------------
	
	\subsection{8. All-Cash Real Estate Transactions}
	A luxury condominium developer in Miami accepts a \$4.2 million all-cash purchase from a buyer represented solely by an attorney who declines to identify the true buyer, citing legal privilege. The funds arrive via a series of wire transfers from three separate shell companies registered in the British Virgin Islands. What is the primary AML red flag?
	\begin{enumerate}[label=\Alph*.]
		\item The developer is in Florida, a historically low-risk real estate market.
		\item The combination of all-cash payment, anonymous shell company ownership, refusal to identify the beneficial owner, and offshore fund sourcing constitutes a classic real estate layering pattern.
		\item The transaction is exempt because a licensed attorney is involved.
		\item The amount is below the \$5 million threshold for mandatory geographic targeting order (GTO) reporting.
	\end{enumerate}
	\textbf{Correct Answer: B}
	
	\subsection{9. FinCEN Geographic Targeting Orders (GTOs)}
	FinCEN has issued Geographic Targeting Orders (GTOs) requiring certain U.S. title insurance companies to identify the natural persons behind legal entities that purchase residential real estate above a defined dollar threshold in designated metropolitan areas. What is the primary money laundering risk that GTOs are designed to address?
	\begin{enumerate}[label=\Alph*.]
		\item Mortgage fraud through falsified income documents submitted to regulated lenders.
		\item The use of all-cash, shell company purchases to integrate illicit wealth into real property without triggering standard bank CDD and reporting requirements.
		\item Zoning violations by commercial real estate developers.
		\item Over-reliance on real estate appraisers who are not licensed AML professionals.
	\end{enumerate}
	\textbf{Correct Answer: B}
	
	\subsection{10. Real Estate as an Integration Vehicle}
	A corrupt foreign official launders \$20 million by (1) moving funds through a offshore account, (2) purchasing a luxury Manhattan penthouse via an anonymous LLC, and (3) renting the apartment at market rate and reporting the rental income on legitimate tax returns. What stage of the money laundering cycle does step (3) represent?
	\begin{enumerate}[label=\Alph*.]
		\item Placement
		\item Layering
		\item Integration
		\item Structuring
	\end{enumerate}
	\textbf{Correct Answer: C}
	
	% -------------------------------------------------------
	\section{Hawala and Informal Value Transfer Systems (IVTS)}
	% -------------------------------------------------------
	
	\subsection{11. The Hawala Mechanism}
	A hawala broker in Dubai receives \$50,000 in cash from a client wishing to send the equivalent value to a family member in Karachi, Pakistan. The broker calls his counterpart broker (hawaladar) in Karachi, who delivers the equivalent in Pakistani Rupees to the recipient. No physical currency crosses the border. What fundamental characteristic makes hawala both efficient and a significant AML risk?
	\begin{enumerate}[label=\Alph*.]
		\item The system is regulated by SWIFT and is therefore fully transparent to global banking supervisors.
		\item The settlement occurs through a trust-based obligation network rather than through the formal banking system, leaving no traceable banking record of the underlying transaction or the identities of the parties involved.
		\item Hawala transactions are explicitly exempt from all FATF Recommendations because they serve migrant workers.
		\item The system is only used for amounts below \$1,000 and therefore falls beneath all reporting thresholds globally.
	\end{enumerate}
	\textbf{Correct Answer: B}
	
	\subsection{12. FATF Recommendation 14: Money or Value Transfer Services (MVTS)}
	Under FATF Recommendation 14, countries are required to ensure that natural or legal persons that provide money or value transfer services are:
	\begin{enumerate}[label=\Alph*.]
		\item Completely prohibited from operating unless they are subsidiaries of licensed commercial banks.
		\item Licensed or registered and subject to AML/CFT monitoring, with unlicensed MVTS providers subject to criminal or administrative sanctions.
		\item Exempt from CDD requirements because their transactions are typically small in value.
		\item Required to obtain a FATF ``MVTS Certification'' issued directly by the FATF Secretariat in Paris.
	\end{enumerate}
	\textbf{Correct Answer: B}
	
	% -------------------------------------------------------
	\section{Casino and Gambling AML}
	% -------------------------------------------------------
	
	\subsection{13. Casino AML Obligations Under the BSA}
	Under U.S. Bank Secrecy Act regulations, casinos with annual gross gaming revenues exceeding \$1 million are treated as financial institutions. Which of the following best describes their core AML reporting obligation for cash transactions?
	\begin{enumerate}[label=\Alph*.]
		\item Casinos have no reporting obligations because gambling proceeds are not considered monetary instruments.
		\item Casinos must file Currency Transaction Reports (CTRs) for cash-in or cash-out transactions exceeding \$10,000 in a gaming day and file SARs for suspicious activity exceeding \$5,000.
		\item Casinos must report only chip purchases, not cash payouts, to FinCEN.
		\item Casinos must report all transactions to the IRS under Form 1099-G, with no obligation to FinCEN.
	\end{enumerate}
	\textbf{Correct Answer: B}
	
	\subsection{14. The ``Wash'' Technique in Casinos}
	A money launderer enters a casino with \$200,000 in criminal proceeds, purchases chips for the full amount, gambles minimally, and then cashes out the chips for a casino check payable to their name, accepting a slight gaming loss as a ``laundering fee.'' Which money laundering stage does this activity represent?
	\begin{enumerate}[label=\Alph*.]
		\item Placement only
		\item Layering only
		\item A combined placement and integration technique, where cash is introduced into a regulated environment and exits as a documented casino check with an apparently legitimate gambling origin.
		\item Proliferation financing via recreational gaming assets.
	\end{enumerate}
	\textbf{Correct Answer: C}
	
	% -------------------------------------------------------
	\section{Money Services Businesses (MSBs)}
	% -------------------------------------------------------
	
	\subsection{15. MSB Registration Requirements}
	Under U.S. federal law (31 CFR \S 1022.380), which of the following entities is required to register with FinCEN as a Money Services Business?
	\begin{enumerate}[label=\Alph*.]
		\item A federally insured commercial bank that offers wire transfers as a service.
		\item A check casher that cashes checks totaling more than \$1,000 for any person in any one day, as well as currency dealers, money transmitters, and issuers or sellers of prepaid access.
		\item Any retail business that accepts cash payments for goods valued over \$500.
		\item A foreign central bank operating a USD reserve account at the Federal Reserve.
	\end{enumerate}
	\textbf{Correct Answer: B}
	
	\subsection{16. MSB Agent Risk}
	A large money transmitter has a nationwide agent network of 4,000 retail agents (convenience stores, grocery stores). A compliance audit reveals that one agent location has processed \$3 million in cash transfers in one month -- far exceeding the agent's expected volume based on its neighborhood demographics. What is the primary compliance risk for the money transmitter (principal)?
	\begin{enumerate}[label=\Alph*.]
		\item The transmitter faces no risk because agents operate independently under their own licenses.
		\item The transmitter is responsible for its agent network's compliance; unexplained volume spikes require investigation, agent-level transaction monitoring, and potential SAR filing by the principal on the agent's activity.
		\item The transmitter must immediately notify local police before conducting any internal review.
		\item The transmitter should increase the agent's transaction limit to avoid customer inconvenience.
	\end{enumerate}
	\textbf{Correct Answer: B}
	
	% -------------------------------------------------------
	\section{Risk-Based Approach (RBA) Framework}
	% -------------------------------------------------------
	
	\subsection{17. Core Principle of the Risk-Based Approach}
	The FATF's risk-based approach (RBA) to AML/CFT compliance allows financial institutions to differentiate the intensity of their due diligence measures. What is the fundamental justification for applying a risk-based approach rather than a purely rules-based, ``one-size-fits-all'' approach?
	\begin{enumerate}[label=\Alph*.]
		\item Regulators prefer rules-based approaches because they are easier to audit.
		\item An RBA concentrates resources and controls on areas of highest risk, making the overall system more effective and efficient, while proportionate measures are applied to lower-risk situations without creating unnecessary burdens.
		\item An RBA eliminates all due diligence requirements for low-risk customers.
		\item The RBA is only applicable to correspondent banking relationships, not retail customers.
	\end{enumerate}
	\textbf{Correct Answer: B}
	
	\subsection{18. Enterprise-Wide Risk Assessment (EWRA)}
	A bank's Enterprise-Wide Risk Assessment (EWRA) must document which of the following?
	\begin{enumerate}[label=\Alph*.]
		\item Only the sanctions screening procedures for correspondent accounts.
		\item An assessment of the institution's inherent risk exposure across all products, services, customer types, geographies, and delivery channels, as well as the mitigating controls in place, resulting in a net residual risk determination.
		\item A list of all suspicious activity reports filed in the previous calendar year sorted by dollar amount.
		\item The board of directors' personal tax filings to assess insider risk.
	\end{enumerate}
	\textbf{Correct Answer: B}
	
	% -------------------------------------------------------
	\section{De-Risking and Financial Inclusion}
	% -------------------------------------------------------
	
	\subsection{19. Consequences of Indiscriminate De-Risking}
	A major international bank terminates its correspondent banking relationships with all respondent banks headquartered in Sub-Saharan Africa, citing ``high risk'' of the region without conducting individual assessments, in order to reduce its AML compliance burden. What is the most significant systemic problem caused by this indiscriminate de-risking approach?
	\begin{enumerate}[label=\Alph*.]
		\item The respondent banks will simply switch to a different SWIFT messaging code.
		\item Legitimate businesses, remittance corridors, and financial institutions in developing economies lose access to the global financial system, driving transactions into unregulated, informal channels that are far less transparent and harder to monitor.
		\item The correspondent bank will automatically receive a reduced FATF mutual evaluation score.
		\item The respondent banks will merge with local central banks, which is the desired regulatory outcome.
	\end{enumerate}
	\textbf{Correct Answer: B}
	
	\subsection{20. FATF and the Wolfsberg Group Guidance on De-Risking}
	The FATF and Wolfsberg Group have both issued guidance addressing the de-risking phenomenon. Which of the following statements is most consistent with that guidance?
	\begin{enumerate}[label=\Alph*.]
		\item Broad, categorical termination of entire classes of customer or geographic relationships is the preferred risk management strategy under the RBA.
		\item Financial institutions should conduct individual, documented risk assessments of respondent institutions rather than applying blanket exclusions based solely on geography or customer type; the RBA does not require zero-risk tolerance.
		\item Regulators require a zero-risk tolerance policy for all high-risk jurisdictions, making de-risking mandatory.
		\item De-risking is only a concern for the correspondent banks; respondent institutions bear no resulting compliance obligations.
	\end{enumerate}
	\textbf{Correct Answer: B}
	
	% -------------------------------------------------------
	\section{Emerging Typologies: NFTs and Online Gaming}
	% -------------------------------------------------------
	
	\subsection{21. Non-Fungible Tokens (NFTs) as a Laundering Vehicle}
	A high-value NFT (valued at \$2 million) is sold between two digital wallets that, on blockchain analysis, are traced back to the same controlling entity. The transaction is recorded publicly on the blockchain at \$2 million. What money laundering technique does this most closely mirror?
	\begin{enumerate}[label=\Alph*.]
		\item Structuring through fractional NFT purchases.
		\item Wash trading / self-dealing, used to artificially establish a price record and move value between controlled wallets while creating the appearance of a legitimate high-value arms-length sale.
		\item Hawala-style informal value transfer.
		\item A standard primary market art auction.
	\end{enumerate}
	\textbf{Correct Answer: B}
	
	\subsection{22. Online Gaming and Virtual Currency Laundering}
	A money launderer purchases large quantities of in-game currency in an online multiplayer game using criminal proceeds through a third-party ``gold-selling'' marketplace. The in-game currency is then transferred to accomplices' accounts, who sell it back on the same marketplace to other players for real-world fiat currency deposited into clean bank accounts. Which stage of the laundering cycle is represented by the final fiat currency recovery?
	\begin{enumerate}[label=\Alph*.]
		\item Placement -- converting cash into the gaming platform.
		\item Layering -- transferring in-game currency between accounts.
		\item Integration -- converting in-game proceeds back into clean fiat currency through peer-to-peer marketplace sales.
		\item Structuring -- breaking the original cash into small gaming transactions.
	\end{enumerate}
	\textbf{Correct Answer: C}
	
	% -------------------------------------------------------
	\section{Predicate Offenses and Proceeds of Crime}
	% -------------------------------------------------------
	
	\subsection{23. FATF Recommendation 3: Predicate Offenses}
	Under FATF Recommendation 3, money laundering should be criminalized based on the widest range of predicate offenses. Which of the following is the FATF's preferred approach?
	\begin{enumerate}[label=\Alph*.]
		\item The all-crimes approach, under which money laundering applies to proceeds from all criminal acts defined as serious under national law.
		\item A narrow list of exactly twelve designated predicate offenses that every country must adopt verbatim.
		\item Only drug trafficking and terrorism-related offenses qualify as predicate crimes under international standards.
		\item Predicate offense lists are entirely optional; each country may choose not to designate any predicate offenses.
	\end{enumerate}
	\textbf{Correct Answer: A}
	
	\subsection{24. Tracing the Proceeds: Comingling}
	A drug trafficker deposits \$300,000 in narcotics proceeds into a legitimate restaurant business, combining them with \$700,000 in genuine food and beverage revenues in a single bank account. The total \$1,000,000 is then used to purchase commercial real estate. What AML challenge does comingling of funds create for investigators?
	\begin{enumerate}[label=\Alph*.]
		\item No challenge -- investigators can simply charge the restaurant owner with tax evasion, which is easier to prove.
		\item Comingling obscures the proportion of illicit funds within the legitimate business account, complicating asset tracing and potentially protecting a portion of the real estate purchase from forfeiture if investigators cannot identify and segregate the criminal proceeds.
		\item Comingling is a legal accounting technique that automatically purifies any criminal proceeds mixed with legitimate revenues.
		\item Investigators need only seize the entire \$1 million, so comingling has no strategic value to the launderer.
	\end{enumerate}
	\textbf{Correct Answer: B}
	
	\subsection{25. Self-Laundering vs. Third-Party Laundering}
	Under most AML frameworks, including FATF standards, can a person who committed the predicate offense be charged with money laundering for processing the proceeds of their own crime?
	\begin{enumerate}[label=\Alph*.]
		\item No -- international standards require that the money launderer and the predicate offender be different individuals.
		\item Yes -- FATF Recommendation 3 explicitly states that the offense of money laundering should apply to persons who committed the predicate offense, i.e., self-laundering is a separate criminal offense.
		\item Only if the predicate offense was committed in a different jurisdiction.
		\item Only if the proceeds exceed \$1 million in aggregate value.
	\end{enumerate}
	\textbf{Correct Answer: B}
	
	% -------------------------------------------------------
	\section{Sanctions: OFAC Special Topics}
	% -------------------------------------------------------
	
	\subsection{26. OFAC General vs. Specific Licenses}
	A U.S. company wishes to engage in a transaction that would normally be prohibited by OFAC sanctions. The company determines there is no applicable General License. What is the next appropriate step?
	\begin{enumerate}[label=\Alph*.]
		\item Proceed with the transaction because the absence of a General License means the transaction must be permissible.
		\item Apply to OFAC for a Specific License, which grants authorization for the particular transaction if OFAC determines it is consistent with U.S. foreign policy objectives.
		\item Route the transaction through a third country to avoid OFAC jurisdiction.
		\item File a SAR instead of a license application to obtain regulatory clearance.
	\end{enumerate}
	\textbf{Correct Answer: B}
	
	\subsection{27. Secondary Sanctions Exposure}
	A European bank processes a EUR-denominated trade finance transaction for a non-U.S. client conducting business with an Iranian entity. The transaction does not touch the U.S. financial system and does not involve USD. What U.S. sanctions risk, if any, does the European bank face?
	\begin{enumerate}[label=\Alph*.]
		\item No risk -- U.S. OFAC jurisdiction is limited to USD transactions and U.S. persons.
		\item Secondary sanctions risk -- the U.S. may impose penalties on the European bank for dealing with a sanctioned jurisdiction even if no U.S. nexus exists, restricting the bank's ability to access U.S. markets and correspondent banking relationships.
		\item The bank is only at risk if the Iranian entity is on the EU sanctions list.
		\item The bank is protected by the EU Blocking Statute, which eliminates all secondary sanctions exposure.
	\end{enumerate}
	\textbf{Correct Answer: B}
	
	% -------------------------------------------------------
	\section{High-Difficulty Scenario: The Layered Real Estate Scheme}
	% -------------------------------------------------------
	
	\noindent A narcotics trafficking network generates \$10 million in drug proceeds in the United States. The proceeds are:
	
	\begin{enumerate}
		\item Deposited in \$9,000 increments into 15 different accounts held by recruited intermediaries across five states.
		\item Aggregated and wired to a Delaware LLC whose only asset is a series of invoices from a fictitious ``consulting'' company incorporated in the Marshall Islands.
		\item Used by the Delaware LLC to purchase a commercial hotel property in Las Vegas for \$9.5 million, with the remaining funds wired to a Cayman Islands bank account.
		\item The hotel operates for 18 months, generating \$2 million in legitimate rental and service revenues. The criminal organization then sells the hotel to an arms-length buyer for \$12 million.
	\end{enumerate}
	
	\subsection{28. Which step constitutes the ``placement'' stage?}
	\begin{enumerate}[label=\Alph*.]
		\item The sale of the hotel to the arms-length buyer.
		\item The purchase of the hotel property by the Delaware LLC.
		\item The structured deposits of \$9,000 increments into the 15 intermediary accounts.
		\item The operation of the hotel and generation of rental revenue.
	\end{enumerate}
	\textbf{Correct Answer: C}
	
	\subsection{29. What financial crime is specifically committed in step (1) of the scheme?}
	\begin{enumerate}[label=\Alph*.]
		\item Wire fraud
		\item Structuring (smurfing) -- deliberately breaking cash deposits into sub-threshold amounts to evade Currency Transaction Report (CTR) requirements, which is a federal crime independent of the underlying predicate offense.
		\item Check kiting
		\item Mortgage fraud
	\end{enumerate}
	\textbf{Correct Answer: B}
	
	\subsection{30. In step (4), what financial red flag pattern should alert the acquiring bank's AML team when the hotel is sold?}
	\begin{enumerate}[label=\Alph*.]
		\item The sale is to an arms-length buyer, so there is no red flag.
		\item The rapid appreciation of property value (\$9.5M purchase to \$12M sale in 18 months), combined with original purchase by an anonymous LLC, wire proceeds to a secrecy-haven Cayman account, and no verified legitimate source of the original acquisition funds -- all constitute integration-stage red flags warranting EDD and a potential SAR.
		\item The only obligation is to file a CTR for the \$12M wire.
		\item The hotel sale is outside the bank's monitoring scope since the bank is not a party to the real estate transaction.
	\end{enumerate}
	\textbf{Correct Answer: B}
	
	
	\section{Advanced CAMS Missing Topics – Expert Level}
	
	\subsection{1. FinCEN 314(a) Information Sharing}
	
	A financial institution receives a FinCEN 314(a) request identifying several individuals suspected of laundering proceeds from transnational organized crime. Which statement is most accurate?
	
	\begin{enumerate}[label=\Alph*.]
		\item The institution may disclose the request to the customer if no match is found.
		\item The institution must search only current customers.
		\item The institution must search accounts and transactions according to FinCEN instructions and maintain confidentiality of the request.
		\item The institution must automatically freeze any matching account.
	\end{enumerate}
	
	\textbf{Correct Answer: C}
	
	\subsection{2. FinCEN 314(b)}
	
	What is the primary purpose of FinCEN 314(b)?
	
	\begin{enumerate}[label=\Alph*.]
		\item Mandatory SAR filing.
		\item Voluntary information sharing among financial institutions regarding suspected money laundering or terrorist financing.
		\item International sanctions enforcement.
		\item Beneficial ownership verification.
	\end{enumerate}
	
	\textbf{Correct Answer: B}
	
	\subsection{3. Corporate Transparency Act}
	
	Which entity is primarily responsible for reporting beneficial ownership information under the Corporate Transparency Act?
	
	\begin{enumerate}[label=\Alph*.]
		\item The reporting company.
		\item The company's primary bank.
		\item The external auditor.
		\item The state tax authority.
	\end{enumerate}
	
	\textbf{Correct Answer: A}
	
	\subsection{4. Source of Wealth vs Source of Funds}
	
	A politically exposed person deposits \$15 million from the sale of a shipping company. Which verification activity addresses Source of Wealth?
	
	\begin{enumerate}[label=\Alph*.]
		\item Reviewing the wire transfer message.
		\item Verifying ownership and historical value creation of the shipping company.
		\item Confirming account number accuracy.
		\item Performing sanctions screening only.
	\end{enumerate}
	
	\textbf{Correct Answer: B}
	
	\subsection{5. Geographic Targeting Orders (GTOs)}
	
	What is the principal objective of FinCEN Geographic Targeting Orders in real estate transactions?
	
	\begin{enumerate}[label=\Alph*.]
		\item Regulate mortgage rates.
		\item Identify beneficial owners behind certain high-value real estate purchases.
		\item Monitor property tax payments.
		\item Reduce closing costs.
	\end{enumerate}
	
	\textbf{Correct Answer: B}
	
	\subsection{6. Egmont Group}
	
	The Egmont Group primarily facilitates:
	
	\begin{enumerate}[label=\Alph*.]
		\item International securities trading.
		\item Information sharing among Financial Intelligence Units.
		\item Global tax collection.
		\item Central bank supervision.
	\end{enumerate}
	
	\textbf{Correct Answer: B}
	
	\subsection{7. Transaction Monitoring Model Validation}
	
	During validation of a transaction monitoring system, investigators discover that 80\% of historical SAR cases would not have generated alerts under current thresholds. What is the most significant concern?
	
	\begin{enumerate}[label=\Alph*.]
		\item Excessive false positives.
		\item Ineffective scenario calibration causing high false negatives.
		\item Regulatory over-reporting.
		\item Inadequate sanctions screening.
	\end{enumerate}
	
	\textbf{Correct Answer: B}
	
	\subsection{8. Data Quality Risk}
	
	Which data issue creates the greatest AML monitoring risk?
	
	\begin{enumerate}[label=\Alph*.]
		\item Duplicate branch codes.
		\item Missing customer identification numbers preventing entity linkage.
		\item Delayed statement generation.
		\item Inconsistent check numbering.
	\end{enumerate}
	
	\textbf{Correct Answer: B}
	
	\subsection{9. Board Oversight}
	
	Which responsibility belongs directly to a financial institution's Board of Directors?
	
	\begin{enumerate}[label=\Alph*.]
		\item Reviewing every alert.
		\item Filing SARs.
		\item Providing oversight of AML risk governance and compliance effectiveness.
		\item Conducting KYC interviews.
	\end{enumerate}
	
	\textbf{Correct Answer: C}
	
	\subsection{10. Artificial Intelligence in AML}
	
	A machine learning monitoring model identifies suspicious activity patterns with high accuracy but cannot explain its decisions. What is the primary regulatory concern?
	
	\begin{enumerate}[label=\Alph*.]
		\item Excessive storage requirements.
		\item Lack of model explainability and governance.
		\item High processing speed.
		\item Reduced investigator workload.
	\end{enumerate}
	
	\textbf{Correct Answer: B}
	
	\subsection{11. Adverse Media Screening}
	
	A customer is not sanctioned and has no criminal convictions but is repeatedly linked through credible media sources to corruption allegations. What is the most appropriate action?
	
	\begin{enumerate}[label=\Alph*.]
		\item Ignore the reports.
		\item Immediately close the account.
		\item Reassess customer risk and consider enhanced due diligence.
		\item File a CTR.
	\end{enumerate}
	
	\textbf{Correct Answer: C}
	
	\subsection{12. Public-Private Partnerships}
	
	Public-private AML partnerships primarily seek to:
	
	\begin{enumerate}[label=\Alph*.]
		\item Replace SAR filing obligations.
		\item Improve intelligence sharing and detection of complex financial crime networks.
		\item Eliminate KYC requirements.
		\item Centralize all investigations within law enforcement.
	\end{enumerate}
	
	\textbf{Correct Answer: B}
	
	\subsection{13. Private Banking Risk}
	
	Which characteristic most increases AML risk within private banking?
	
	\begin{enumerate}[label=\Alph*.]
		\item Small account balances.
		\item Limited international exposure.
		\item Complex ownership structures and large cross-border transactions.
		\item Domestic payroll deposits.
	\end{enumerate}
	
	\textbf{Correct Answer: C}
	
	\subsection{14. GDPR and AML Requirements}
	
	A European bank receives a law enforcement request related to a money laundering investigation. Which principle generally applies?
	
	\begin{enumerate}[label=\Alph*.]
		\item AML obligations automatically cease under GDPR.
		\item Data protection and AML requirements must be balanced according to applicable legal exemptions.
		\item Customer consent is always required.
		\item Information cannot be shared under any circumstances.
	\end{enumerate}
	
	\textbf{Correct Answer: B}
	
	\subsection{15. Regulatory Root Cause Analysis}
	
	Following a regulatory enforcement action, investigators identify weak governance, poor data quality, and inadequate staffing. These findings collectively represent:
	
	\begin{enumerate}[label=\Alph*.]
		\item Isolated operational incidents.
		\item Root causes requiring enterprise-wide remediation.
		\item Customer onboarding errors only.
		\item Cybersecurity deficiencies exclusively.
	\end{enumerate}
	
	\textbf{Correct Answer: B}
\end{document}